\section{Objectives and the question}
\label{sec:intro}
\label{sec:objectives}
\textbf{Purpose.} ARGUS asks how the outer wing of a transport aircraft should adapt during
cruise. A low-order VSPAERO study searched that design space and selected the candidate
shapes~\cite{zheng}. The search is inviscid, wing-only, and ranks on induced drag. This work is
the viscous audit of that search.

\textbf{Study hierarchy.} The design study evaluates hundreds of geometries cheaply. This
campaign evaluates six of them with steady RANS at the matched Reynolds number. The design study
proposes candidates. The RANS says what they cost at trimmed lift.

\textbf{The job.} Six wing surfaces, one rigid baseline and five morphing candidates, were meshed and
solved with steady RANS at three operating points, every geometry trimmed to that point's target
lift. The low-order study selected the candidates at the low-speed point, Condition~CR, so that is
where it is tested against the RANS; the two cruise points are the results the low-order study
does not have.
\textbf{Question.} This report answers two. First, does the low-order method order the candidates
the way the RANS orders them? Second, what does the RANS say about the candidates themselves?

\textbf{Reading order.} The report follows the order of the work. Section~\ref{sec:condgeom} states
the operating points and the six geometries. Section~\ref{sec:method} states the OpenFOAM setup,
the automated pipeline, the trim convention and the convergence gate, each once.
Section~\ref{sec:qoi} lists every quantity of interest, its frame, and whether the low-order study
produces it. Section~\ref{sec:results} reports each quantity in turn, RANS beside VLM wherever a
VLM value exists. Section~\ref{sec:lo} answers whether the VLM can drive the design study, and
Section~\ref{sec:concl} states what the whole of it establishes and what it does not.
Appendices~\ref{sec:app:sections} and~\ref{sec:deltamaps} carry the surface pressure and skin
friction, by station and over the whole wing.
